%\clearpage

\section{Summary of changes included in version 9 of T2K-TN-141 (released on Sept. 11, 2013)}

This section is designed to allow readers to keep track of the changes between versions of T2K-TN-141.
Table, figure, equation, footnote or citation numbers shown in black refer to version of 8,
{\color{red}whereas the numbers shown in red refer to version of 9}.
Both numbers are shown where possible.

\begin{itemize}

\item
  Information on, and studies involving fitting for $\sin^2\theta_{23}$, which was first shown
  in Fig.\ 90 {\color{red}101} and Tab.\ 14 ({\color{red}13}).
  All additional fits have been done with the
  '2+N' fitter, and use the Run 1+2+3 dataset only.
  \begin{itemize}
  \item The reasoning for fitting for $\sin^2\theta_{23}$ has been added to Sec.\ 4.1.1 {\color{red}(4.1.1)}.
  \item A footnote has been added to Sec.\ 4.3 {\color{red}(4.3)} describing fitter performance tests.
  \item The best-fit point has been added to Sec.\ 5.1 {\color{red}(5.1)}.
  \item The p-value has been added to Tab.\ 14 {\color{red}(13)}.
  \item The pulls of the systematics at the best-fit point is shown in Fig.\ {\color{red}33}.
  \item The goodness of fit is shown in Fig.\ {\color{red}41}.
  \item The best-fit spectrum is shown in Fig.\ {\color{red}45} for the fine binning scheme,
    Fig.\ {\color{red}53} for the coarse binning scheme, and Fig.\ {\color{red}54} for a breakdown by MC component prediction.
  \item A comparison of the best-fit spectra of the '2+N' fit for $\sin^2\theta_{23}$
    and the fit for $\sin^22\theta_{23}$ assuming the second octant is shown in Figs.\ {\color{red}61} and {\color{red}62},
    for the fine and coarse binning schemes respectively.
  \item The 1D constant $\Delta \chi^2$ 68\% and 90\% confidence limits have been added to Section 5.4 {\color{red}(5.4)}.
    The 1D profile likelihoods are shown in Figs.\ {\color{red}67} and {\color{red}68}
    for $\sin^2\theta_{23}$ and $|\Delta m^2_{32}|$ respectively.
  \item The $\Delta \chi^2$ surface is shown in Fig.\ {\color{red}71}.
  \item The fractional difference between $\Delta \chi^2$ surfaces for both octant assumptions, are shown in Fig.\ {\color{red}73}.
  \item The contour made using the constant $\Delta \chi^2$ method is shown in Fig.\ {\color{red}88}.
  \item A comparison between the $\sin^2\theta_{23}$ contour with external data is shown in Fig.\ {\color{red}108}.
  \item The result of the '2+N' $\sin^2\theta_{23}$ best-fit point and the corresponding 1D 90\% C.L. limits have been added
    to both the Abstract, and Sec.\ 6 {\color{red}(6)}.
  \end{itemize}

\end{itemize}

